\section*{Chapter \ref{ch:s1:seasonality-and-calendar-effects} --- Seasonality and Calendar Effects}

\begin{solution}{exo:s1:seasonality-and-calendar-effects:1}
$100 \times 0.05 = 5$; with 75 null rules, $75 \times 0.05 = 3.75$. The experiment found 8, within the noise of correlated rules (a false weekday
tends to bring false days and fortnights with it).
\end{solution}

\begin{solution}{exo:s1:seasonality-and-calendar-effects:2}
$\Phi^{-1}(1 - 0.05/200) = 3.48$.
\end{solution}

\begin{solution}{exo:s1:seasonality-and-calendar-effects:3}
Nobody can buy January gas in July and hold it as spot; the tradeable instruments are futures, whose prices for winter delivery already include the
expected winter premium. A futures position earns only the difference between the realised and the priced season, plus any risk premium.
\end{solution}

\begin{solution}{exo:s1:seasonality-and-calendar-effects:4}
Daily volatility is $0.147/\sqrt{252} = 0.93\%$. For the turn of the month the standard error is $0.93\% \times \sqrt{1/1{,}440 + 1/6{,}120} = 2.71$ basis
points, so 8 basis points give $t = 2.95$ on average (the sample drawn gave 4.4). For holidays it is $0.93\% \times \sqrt{1/270 + 1/7{,}290} = 5.74$
basis points and $t = 2.61$ (the sample gave 1.5).
\end{solution}

\begin{solution}{exo:s1:seasonality-and-calendar-effects:5}
The t-statistic grows with the square root of the sample: $30 \times (3.48/2.61)^2 = 53$ years.
\end{solution}

\begin{solution}{exo:s1:seasonality-and-calendar-effects:6}
It controls the expected share of false discoveries among those made, not the chance of making any; with many true effects it keeps more of them and
tolerates a few false ones. That is acceptable when discoveries feed further tests (a research pipeline), not when each is traded as found.
\end{solution}

\begin{solution}{exo:s1:seasonality-and-calendar-effects:7}
With seed 8 the best turn-of-month rule has $t = 2.25$ and the pre-holiday rule $t = 2.99$: uncorrected, 8 turn-of-month rules, the pre-holiday rule
and 4 rules that carry nothing are found; every correction finds nothing. The planted effects did not change; the sample did. What should stay is the
pattern of the corrections (few or no false discoveries); what changes is which true effects are lucky enough to pass.
\end{solution}

\begin{solution}{exo:s1:seasonality-and-calendar-effects:8}
Mondays are one of many calendar rules anyone has looked at; $t = 2.4$ does not survive a correction for them. A mechanism, an out-of-sample period and
the costs of 52 round trips a year are needed before trading it.
\end{solution}

\begin{solution}{pb:s1:seasonality-and-calendar-effects:1}
\begin{enumerate}
\item A return difference on calendar-defined days; higher returns on the last and first days of the month; recurrence of relative returns at the same
point of each year.
\item $+8.8\%$ in January, $+7.3\%$ in December, $-5.3\%$ in March and $-4.2\%$ in April against trend: winter heating, spring refilling.
\item October and November lowest ($t = -1.71$, $-2.17$), nothing after correction (smallest adjusted p-value 0.36).
\item Persistently anomalous returns around the turn of the week, month and year and around holidays, over ninety years.
\item Twelve 21-day months, nine holidays, thirty years; 8 basis points on four turn-of-month days, 15 on pre-holiday days; a hundred rules.
\item By their expected difference from each planted effect: at least 3 basis points carries it.
\item None: 19 turn-of-month, 0 pre-holiday, 8 false; Bonferroni, Holm and Benjamini--Yekutieli: 6, 0, 0; Benjamini--Hochberg: 11, 0, 3.
\item Nine days a year give 270 observations in thirty years: an expected t of 2.6, 1.5 in this sample.
\item Stocks' relative returns recur in the same calendar month for up to twenty annual lags.
\item A 1\% monthly seasonal return per stock; the average return in the same month of the previous five years.
\item A rank IC of 0.031 ($t = 3.8$); a decile spread of 0.50\% a month, Sharpe ratio 0.72 before costs.
\item Each stock shows it once a year, against a monthly noise of several percent.
\item \textbf{Named result.} The \omterm{def:s1:seasonality-and-calendar-effects:tom}{turn-of-the-month effect} is found under every correction (6 rules under the family-wise ones, 11 under
Benjamini--Hochberg), the pre-holiday effect under none; false discoveries number 8 uncorrected, 3 under Benjamini--Hochberg and 0 under Bonferroni,
Holm and Benjamini--Yekutieli.
\item A mechanism stated first, a test that counts every rule tried, and enough observations.
\item A family-wise correction (Holm) for a search to be traded directly; the false discovery rate for a screen followed by more tests.
\item The window fixed in advance, every window tried reported, futures prices and costs, an out-of-sample period.
\item That testing many rules on the same data finds effects that are not there.
\item The pre-holiday effect.
\item Book~4 derived the corrections; here they decide which \omterm{def:s1:seasonality-and-calendar-effects:calendar}{calendar effects} are believed.
\item Those with a mechanism that survive a correction for every rule tried.
\end{enumerate}
\end{solution}

\begin{solution}{iq:s1:seasonality-and-calendar-effects:1}
Count how many calendar rules have been tried (by me and by the literature), correct for them, look for a mechanism, check other markets and periods,
and estimate what it would earn after costs; most likely, file it.
\end{solution}

\begin{solution}{iq:s1:seasonality-and-calendar-effects:2}
The family-wise error rate is the probability of at least one false discovery; the false discovery rate is the expected share of discoveries that are
false. The first is stricter; the second keeps more true effects when there are many.
\end{solution}

\begin{solution}{iq:s1:seasonality-and-calendar-effects:3}
Salaries, pension contributions and fund inflows arrive around month-end and are invested within days; \omterm{def:s1:index-rebalancing:fund}{index funds} rebalance at month-end. Check fund
flow data by day of month and whether the effect is larger where such flows are larger.
\end{solution}

\begin{solution}{iq:s1:seasonality-and-calendar-effects:4}
Each exchange's trading days and holidays, as they were each year; early closes; expiry and settlement dates; time zones; and the mapping of events
known in advance from those announced later.
\end{solution}

\begin{solution}{iq:s1:seasonality-and-calendar-effects:5}
Spot seasonality is predictable and priced into the futures curve, so futures returns need not be seasonal; seasonality in futures returns is a risk
premium or a mispricing, and much harder to find.
\end{solution}

\begin{solution}{iq:s1:seasonality-and-calendar-effects:6}
The difference of means has standard error $\sigma\sqrt{1/(fn) + 1/((1 - f)n)} = \sigma/\sqrt{f(1 - f)n}$, so $E[t] = \delta\sqrt{f(1 - f)n}/\sigma$; it
reaches $z$ when $n = z^2\sigma^2/(\delta^2 f(1 - f))$.
\end{solution}
